
\documentclass{article}
\usepackage{colortbl}
\usepackage{makecell}
\usepackage{multirow}
\usepackage{supertabular}

\begin{document}

\newcounter{utterance}

\centering \large Interaction Transcript for game `matchit\_ascii', experiment `similar\_grid\_2', episode 7 with Qwen3.6{-}35B{-}A3B{-}FP8{-}without{-}reasoning.
\vspace{24pt}

{ \footnotesize  \setcounter{utterance}{1}
\setlength{\tabcolsep}{0pt}
\begin{supertabular}{c@{$\;$}|p{.15\linewidth}@{}p{.15\linewidth}p{.15\linewidth}p{.15\linewidth}p{.15\linewidth}p{.15\linewidth}}
    \# & $\;$A & \multicolumn{4}{c}{Game Master} & $\;\:$B\\
    \hline

    \theutterance \stepcounter{utterance}  
    & & \multicolumn{4}{p{0.6\linewidth}}{
        \cellcolor[rgb]{0.9,0.9,0.9}{
            \makecell[{{p{\linewidth}}}]{
                \texttt{\tiny{[P1$\langle$GM]}}
                \texttt{İşbirlikçi bir tahmin oyununa katılıyorsun.} \\
\texttt{Amaç, bu ızgara ile yalnızca benim görebildiğim diğer görselin aynı olup olmadığını anlamak.} \\
\texttt{Önce sen kendi ızgaranı tarif et. Ardından ben de kendi görselimin tarifini vereceğim ve ızgaraların aynı olup olduğunu anlamak için birbirimize sorular sorabileceğiz.  Şimdi ızgaranın kısa bir tarifini "\textbackslash{}"AÇIKLAMA:\textbackslash{}"" ile başlayacak şekilde yaz ve ardından açıklamanı ekle. Başka hiçbir şey ekleme. İşte senin ızgaran:} \\
\texttt{▢ ▢ X ▢ ▢} \\
\texttt{▢ ▢ X ▢ ▢} \\
\texttt{▢ ▢ X ▢ ▢} \\
\texttt{▢ X ▢ X ▢} \\
\texttt{X ▢ ▢ ▢ X} \\
            }
        }
    }
    & \\ \\

    \theutterance \stepcounter{utterance}  
    & \multicolumn{4}{p{0.6\linewidth}}{
        \cellcolor[rgb]{0.9,0.9,0.9}{
            \makecell[{{p{\linewidth}}}]{
                \texttt{\tiny{[P1$\rangle$GM]}}
                \texttt{AÇIKLAMA: 5x5 boyutunda bir ızgara. Üçüncü sütun (soldan) tamamen X ile dolu. Dördüncü satırda ikinci ve dördüncü hücreler X. Beşinci satırda birinci ve beşinci hücreler X. Geri kalan tüm hücreler boş.} \\
            }
        }
    }
    & & \\ \\

    \theutterance \stepcounter{utterance}  
    & & & \multicolumn{2}{p{0.3\linewidth}}{
        \cellcolor[rgb]{0.9,0.9,0.9}{
            \makecell[{{p{\linewidth}}}]{
                \texttt{\tiny{[GM$|$GM]}}
                \texttt{continue} \\
            }
        }
    }
    & & \\ \\

    \theutterance \stepcounter{utterance}  
    & & \multicolumn{4}{p{0.6\linewidth}}{
        \cellcolor[rgb]{0.9,0.9,0.9}{
            \makecell[{{p{\linewidth}}}]{
                \texttt{\tiny{[GM$\rangle$P2]}}
                \texttt{İşbirlikçi bir tahmin oyununa katılıyorsun.} \\
\texttt{Amaç, bu ızgara ile yalnızca benim görebildiğim diğer görselin aynı olup olmadığını anlamak.} \\
\texttt{Önce sen kendi ızgaranı tarif et. Ardından ben de kendi görselimin tarifini vereceğim ve ızgaraların aynı olup olduğunu anlamak için birbirimize sorular sorabileceğiz.  Şimdi ızgaranın kısa bir tarifini "\textbackslash{}"AÇIKLAMA:\textbackslash{}"" ile başlayacak şekilde yaz ve ardından açıklamanı ekle. Başka hiçbir şey ekleme. İşte senin ızgaran:} \\
\texttt{▢ ▢ X ▢ ▢} \\
\texttt{▢ ▢ X ▢ ▢} \\
\texttt{▢ ▢ X ▢ ▢} \\
\texttt{▢ ▢ ▢ X ▢} \\
\texttt{X ▢ X ▢ X} \\
            }
        }
    }
    & \\ \\

    \theutterance \stepcounter{utterance}  
    & & & \multicolumn{4}{p{0.6\linewidth}}{
        \cellcolor[rgb]{0.9,0.9,0.9}{
            \makecell[{{p{\linewidth}}}]{
                \texttt{\tiny{[GM$\langle$P2]}}
                \texttt{AÇLAMA: 5x5'lük bir ızgara. Üstteki üç satırda, üçüncü sütunda 'X' var. Dördüncü satırda, dördüncü sütunda 'X' var. Beşinci satırda ise birinci, üçüncü ve beşinci sütunlarda 'X' var.} \\
            }
        }
    }
     \\ \\

    \theutterance \stepcounter{utterance}  
    & & & \multicolumn{2}{p{0.3\linewidth}}{
        \cellcolor[rgb]{0.9,0.9,0.9}{
            \makecell[{{p{\linewidth}}}]{
                \texttt{\tiny{[GM$|$GM]}}
                \texttt{abort, first word: AÇLAMA:, but expected AÇIKLAMA:} \\
            }
        }
    }
    & & \\ \\

\end{supertabular}
}

\end{document}
